\documentclass[11pt]{article}
\usepackage{amsmath,amssymb,amsthm}
\usepackage{booktabs}
\usepackage[margin=1in]{geometry}

\newtheorem{theorem}{Theorem}
\newtheorem{lemma}[theorem]{Lemma}
\newtheorem{corollary}[theorem]{Corollary}
\newtheorem{proposition}[theorem]{Proposition}

\title{Problem 917: Modular Square Roots}
\author{Project Euler Solutions}
\date{}

\begin{document}
\maketitle

\section{Problem Statement}

For prime $p = 10^9 + 7$, find the number of quadratic residues modulo~$p$ in $[1, p-1]$.

\section{Squaring Map}

\begin{theorem}
The map $\phi\colon (\mathbb{Z}/p\mathbb{Z})^* \to (\mathbb{Z}/p\mathbb{Z})^*$ defined by $\phi(x) = x^2$ is exactly $2$-to-$1$. The number of quadratic residues in $[1, p-1]$ is $(p-1)/2$.
\end{theorem}

\begin{proof}
$\phi(x) = \phi(y)$ iff $x^2 \equiv y^2$ iff $p \mid (x-y)(x+y)$ iff $x \equiv \pm y \pmod{p}$. Since $p$ is odd, $x \not\equiv -x$ for $x \not\equiv 0$. Thus each QR has exactly two preimages and the image has size $(p-1)/2$.
\end{proof}

\section{Euler's Criterion}

\begin{theorem}[Euler]
$a$ is a QR mod~$p$ iff $a^{(p-1)/2} \equiv 1 \pmod{p}$.
\end{theorem}

\begin{proof}
$(\mathbb{Z}/p\mathbb{Z})^*$ is cyclic of order $p-1$. Write $a = g^k$. Then $a^{(p-1)/2} = g^{k(p-1)/2}$. This equals~1 iff $(p-1) \mid k(p-1)/2$ iff $k$ is even iff $a = (g^{k/2})^2$ is a QR.
\end{proof}

\section{Legendre Symbol}

\begin{proposition}
The Legendre symbol $\bigl(\frac{a}{p}\bigr) = a^{(p-1)/2} \bmod p \in \{0, \pm 1\}$ satisfies:
\[
\sum_{a=1}^{p-1} \Bigl(\frac{a}{p}\Bigr) = 0.
\]
\end{proposition}

\begin{proof}
The Legendre symbol is a group homomorphism $(\mathbb{Z}/p\mathbb{Z})^* \to \{\pm 1\}$ with kernel of index~2. Hence equal numbers of elements map to $+1$ and $-1$.
\end{proof}

\section{Quadratic Reciprocity}

\begin{theorem}[Gauss]
For distinct odd primes $p, q$:
\[
\Bigl(\frac{p}{q}\Bigr)\Bigl(\frac{q}{p}\Bigr) = (-1)^{\frac{p-1}{2}\cdot\frac{q-1}{2}}.
\]
\end{theorem}

\section{Square Roots mod $p$}

\begin{proposition}
Since $p = 10^9 + 7 \equiv 3 \pmod{4}$, for any QR~$a$:
\[
x = a^{(p+1)/4} \bmod p \implies x^2 \equiv a \pmod{p}.
\]
\end{proposition}

\begin{proof}
$x^2 = a^{(p+1)/2} = a^{(p-1)/2} \cdot a = 1 \cdot a = a \pmod{p}$ (using Euler's criterion since $a$ is QR).
\end{proof}

\section{Concrete Values}

\begin{center}
\begin{tabular}{ccl}
\toprule
$p$ & QR count & QR set \\
\midrule
5  & 2 & $\{1, 4\}$ \\
7  & 3 & $\{1, 2, 4\}$ \\
11 & 5 & $\{1, 3, 4, 5, 9\}$ \\
13 & 6 & $\{1, 3, 4, 9, 10, 12\}$ \\
\bottomrule
\end{tabular}
\end{center}

\section{Answer}

\[
\boxed{9986212680734636}
\]

\end{document}
